\begin{table}[H]
\centering\centering
\caption{\label{tbl-paper-language-behavior}Linguistic Distance and Language-Learning Behavior. Source: Census of India 2001 bilingualism and trilingualism tables.}
\centering
\fontsize{9}{11}\selectfont
\begin{threeparttable}
\begin{tabular}[t]{>{\raggedright\arraybackslash}p{4.2cm}>{\centering\arraybackslash}p{2.35cm}>{\centering\arraybackslash}p{2.35cm}>{\centering\arraybackslash}p{2.35cm}>{\centering\arraybackslash}p{2.35cm}>{\centering\arraybackslash}p{2.35cm}>{\centering\arraybackslash}p{2.35cm}>{\centering\arraybackslash}p{2.35cm}}
\toprule
\multicolumn{1}{c}{\textbf{ }} & \multicolumn{4}{c}{\textbf{All states}} & \multicolumn{3}{c}{\textbf{Hindi-belt states}} \\
\cmidrule(l{3pt}r{3pt}){2-5} \cmidrule(l{3pt}r{3pt}){6-8}
  & (1) & (2) & (3) & (4) & (5) & (6) & (7)\\
\midrule
\cellcolor{gray!10}{Linguistic distance from Hindi} & \cellcolor{gray!10}{2.234} & \cellcolor{gray!10}{} & \cellcolor{gray!10}{-9.694***} & \cellcolor{gray!10}{-3.581***} & \cellcolor{gray!10}{-0.037} & \cellcolor{gray!10}{} & \cellcolor{gray!10}{-9.024*}\\
 & (1.448) &  & (2.573) & (0.916) & (2.942) &  & (5.321)\\
\cellcolor{gray!10}{Distant-language indicator} & \cellcolor{gray!10}{} & \cellcolor{gray!10}{16.727***} & \cellcolor{gray!10}{} & \cellcolor{gray!10}{} & \cellcolor{gray!10}{} & \cellcolor{gray!10}{10.165} & \cellcolor{gray!10}{}\\
 &  & (4.394) &  &  &  & (10.807) & \\
\cellcolor{gray!10}{Outcome} & \cellcolor{gray!10}{English acquisition} & \cellcolor{gray!10}{English acquisition} & \cellcolor{gray!10}{Hindi acquisition} & \cellcolor{gray!10}{Multilingualism} & \cellcolor{gray!10}{English acquisition} & \cellcolor{gray!10}{English acquisition} & \cellcolor{gray!10}{Hindi acquisition}\\
Population & Multilingual speakers & Multilingual speakers & Multilingual speakers & Native speakers & Multilingual speakers & Multilingual speakers & Multilingual speakers\\
\cellcolor{gray!10}{Observations} & \cellcolor{gray!10}{1,566} & \cellcolor{gray!10}{1,566} & \cellcolor{gray!10}{1,566} & \cellcolor{gray!10}{1,587} & \cellcolor{gray!10}{560} & \cellcolor{gray!10}{560} & \cellcolor{gray!10}{560}\\
Partial $R^2$ & 0.025 & 0.128 & 0.094 & 0.139 & 0.000 & 0.046 & 0.121\\
\cellcolor{gray!10}{State fixed effects} & \cellcolor{gray!10}{Yes} & \cellcolor{gray!10}{Yes} & \cellcolor{gray!10}{Yes} & \cellcolor{gray!10}{Yes} & \cellcolor{gray!10}{Yes} & \cellcolor{gray!10}{Yes} & \cellcolor{gray!10}{Yes}\\
Language-share controls & Yes & Yes & Yes & Yes & Yes & Yes & Yes\\
\bottomrule
\end{tabular}
\begin{tablenotes}[para]
\item Each column is a weighted state-by-language regression. Standard errors in parentheses use HC1 heteroskedasticity-robust inference. All specifications include state fixed effects, native-language state share, an indicator for the state's most common native language, and a Hindi/Urdu reference indicator where required by the registered specification. Continuous-distance coefficients are percentage-point changes per one Shastry distance degree; the distant-language indicator equals one for languages at least three degrees from Hindi. Partial R-squared is the one-degree-of-freedom model-based partial R-squared for the reported linguistic-distance coefficient. Columns (5)-(7) apply the predeclared Hindi-belt sample restriction. These are descriptive associations with language-learning behavior rather than instrumental-variable first stages. * p < 0.10, ** p < 0.05, *** p < 0.01 Source: Census of India 2001 bilingualism and trilingualism tables.
\end{tablenotes}
\end{threeparttable}
\end{table}
